\section{The cycle-profile theorem}

At the chosen visible basepoint $0$, iterating the lifted successor through
one visible circuit gives
\[
T_v^n(0,x)=(0,xh),
\qquad
h=v_0v_1\cdots v_{n-1}.
\]
At a general visible state $i$, define the based circuit holonomy
\[
h_i
=
v_i v_{i+1}\cdots v_{n-1}v_0\cdots v_{i-1}.
\]
Then
\[
T_v^n(i,x)=(i,xh_i).
\]
Every $h_i$ is conjugate to $h$, so all based circuit holonomies have the
same order. Let
\[
d=\ord(h)=\ord(h_i).
\]
Right multiplication by $h_i$ partitions $R$ into $|R|/d$ orbits, each of
length $d$.

\begin{theorem}[Cycle profile]
If the circuit holonomy has order $d$, then the lifted successor has exactly
\[
\frac{|R|}{d}
\]
orbits. Every orbit has length $nd$. Equivalently, the lifted cycle profile is
\[
\underbrace{[nd,\ldots,nd]}_{|R|/d\text{ entries}}.
\]
\end{theorem}

\begin{proof}
A lifted state over visible coordinate $i$ can return to that coordinate only
after a multiple of $n$ steps. After $qn$ steps its receipt coordinate is
\[
xh_i^q.
\]
Because $h_i$ is conjugate to $h$, its order is $d$. Therefore the smallest
positive $q$ returning the receipt coordinate is $d$, and every lifted orbit
has length $nd$. Since the lifted state set has cardinality $n|R|$, the
number of orbits is
\[
\frac{n|R|}{nd}=\frac{|R|}{d}.
\]
\end{proof}

The cycle profile depends only on receipt order, while the finer gauge class
retains the full holonomy conjugacy class.

